\section{Boolean search, equality exploration, and external engines}
\label{sec:boolean}

\subsection{A separate lane before a core-engine rewrite}

Large combinatorial Boolean search is not what a congruence-closure and
theory-propagation engine is designed for, and the Lean reference says as much:
it distinguishes \grind{}'s branching strategy from bulk Boolean solving and
recommends \code{bv\_decide} for the latter~\cite{grind-ref}. These are
complementary roles. Nothing here is evidence that \grind{} is a weak
simplifier, and the existing shared-state \code{grind => bv\_decide} route
already encodes equivalence classes and theory-learned facts into the SAT
problem~\cite{lean-release-434}.

The proposal is therefore a \emph{separate lane} with an explicit proof
contract, scheduled after structural steps have made the relevant facts
available, and implemented late. Every roadmap in this collection places the
Boolean adapter last, and that ordering should be kept even though it is the
component with the most complete prototype.

\subsection{A proof-producing SAT/theory loop}

The coarse-grained version needs no changes to any existing engine:

\begin{lstlisting}
while budget remains:
  answer := SAT.solve(currentClauses)
  if answer is UNSAT:
    replay the propositional certificate and all clause origins
    lift the contradiction to the original Lean goal
    return PROVED only if replay and the axiom audit succeed
  if answer is not a complete Boolean assignment:
    return UNKNOWN
  open a temporary scope with the assignment's signed atoms
  ask the relevant Lean theory engines for a contradiction
  if a contradiction is proved:
    discharge its used decisions into a learned clause
    add the proved clause; close the temporary scope
  else:
    refine the supported abstraction or return UNKNOWN
\end{lstlisting}

Incremental CDCL(T) is the later optimisation. When a theory implies a literal
$L$ from current literals $A_1,\ldots,A_k$, it must return a proof of the clause
$\neg A_1\lor\cdots\lor\neg A_k\lor L$. A stale explanation referring to a
withdrawn equality is a correctness bug even when the Boolean clause still looks
plausible.

\subsection{An implemented CDCL(T) with checkable learning}
\status{Implemented and tested in Python by two independent efforts, with two
different proof formats.}

The implemented theory is integer difference logic: atoms $x-y\le c$ over
integer variables, with Boolean structure in CNF. The solver performs unit
propagation to a Boolean fixed point, then tests the implied difference
constraints for a negative cycle; a negative cycle yields a theory lemma. First-UIP
conflict analysis and non-chronological backjumping are implemented, and
tautological resolvents are detected during analysis.

\paragraph{The integer-negation trap.}
The negation of $x-y\le c$ over $\Z$ is $x-y\ge c+1$, that is $y-x\le -c-1$ ---
\emph{not} $y-x\le -c$. Getting this wrong produces a solver that is subtly
unsound on exactly the instances it is meant to decide. Two defences are worth
copying. The consistency oracle carries a dedicated regression assertion at the
complement boundary, tested at magnitudes such as $10^{60}$ to confirm that
arbitrary-precision bounds survive; and the certificate checker \emph{duplicates}
the negation rule rather than calling the solver's shared edge constructor, so
that a bug in one copy cannot cancel itself out in the other. For the same
reason the all-pairs consistency check uses explicit \code{None} sentinels for
absent edges instead of a finite infinity.

\paragraph{Two certificate formats.}
One implementation emits a \emph{resolution DAG}: nodes tagged
\code{input}, \code{theory}, or \code{resolution}, each resolution node carrying
parent indices and a pivot. Replay is purely structural and cheap --- it re-does
the solver's own steps, which makes it fast but means the checker trusts the
solver's choice of derivation.

The other emits \emph{RUP additions plus signed negative-cycle lemmas}, and its
checker re-derives each added clause by \emph{independent} unit propagation. It
never trusts the solver's resolution chain. That independence is the stronger
property, and it is why a merged implementation should offer the resolution DAG
as an optional second mode --- strictly cheaper to check when available --- while
keeping RUP as the default.

A theory lemma is validated by exact coefficient summation: the cycle's
coefficients telescope, and the resulting constant bound must be strictly
negative. A satisfying assignment is validated separately: the Boolean
assignment must satisfy every original clause, \emph{and} each atom's truth
value must match the integer model.

\begin{proposition}[Decision procedure for the implemented fragment]
For a fixed finite set of integer difference atoms and a CNF over them,
exhaustive Boolean search together with a complete difference-logic consistency
check is a decision procedure.
\end{proposition}

This does not extend automatically to the other theories or the quantified goals
\forge{} targets, and it is stated only to fix what the implemented fragment
does.

\paragraph{Production requirements.}
A Lean port needs the clause origins bound to actual Lean propositions, a
supported certificate format, and a declared trust profile. Existing work is
directly relevant: \code{bv\_check} already demonstrates the useful separation
of a stored SAT certificate from an installed solver, and PBLean shows that
existing Lean SAT integrations make an explicit compiled-reflection
choice~\cite{bv-decide-api,pblean}. The classical DPLL(T) framework supplies the
architecture; the contribution proposed here is an integration contract with
native Lean terms, not the algorithm~\cite{dpll-t}.

\begin{remark}[A benchmark fairness rule]
A comparison must not measure strict kernel replay in one tactic against native
replay in another and report them as identical trust configurations. Caching an
LRAT file for \code{bv\_check} does not by itself establish a kernel-only replay
profile.
\end{remark}

\subsection{Bounded equational exploration}

Equality saturation should be used \emph{selectively}, because \grind{} already
supplies congruence machinery. The proposed extension maintains small
collections of alternative forms for terms relevant to a blocked obligation,
prioritising forms that expose recursive calls, match a theorem conclusion, or
reduce an arithmetic certificate's dimension. Rewrites are justified by equality
proofs; conditional rewrites create explicit guards, and if a conditional rule
depends on an unproved proposition, its output belongs in the obligation graph
rather than in an unconditional equality class.

Do not saturate associativity, commutativity, and distributivity
indiscriminately: existing normalisers already handle much of that structure
more efficiently. Use algebraic normal forms for certificate checking while
retaining a few factored or recursive views for search, and limit term count,
expansion degree, and rewrite generations.

Extraction deserves more care than it usually receives. At each fuel boundary,
extract \emph{several} candidate representatives rather than a single
syntactically cheapest term: one may have the fewest nodes, another the simplest
arithmetic degree, another the best overlap with a desired lemma. A practical
cost vector is lexicographic --- unsupported constructs; unsolved side
conditions; maximum polynomial degree or recursive mismatch count; expression
size; estimated replay size --- and different obligation kinds lead with
different coordinates. Cyclic equivalence classes require a well-founded
extraction procedure, because a cycle is not a finite Lean term.

An e-class identifier is meaningful only within a state lineage and revision. A
proof node must not depend on a mutable representative number that does not
survive rollback. The binder-sensitivity and dependent-transport obligations
involved in doing this over real Lean expressions are documented
elsewhere~\cite{egg-lean-paper,eqsat-paper}, and the historical \code{lean-egg}
package is deprecated in favour of \grind{}~\cite{lean-egg}; an archived
experimental project should not become a critical dependency because its name
sounds aligned with the design.

\subsection{Higher-order structure and external provers}

Before translating a goal to a first-order backend, expose its constructive
structure. For a function equality, attempt the relevant extensionality
principle and introduce a fresh argument. For a universally quantified goal,
introduce its binder. For a conjunction or constructor target, create the
corresponding subgoals. For an existential, construct a scoped witness and then
prove its specification.

A bounded type-directed term synthesiser handles witness goals beyond the affine
fragment. Its grammar includes local terms, lambda abstraction, application of
retrieved lemmas and functions, constructors, and selected recursor or fold
schemas, with every candidate elaborated before semantic checking. Higher-order
pattern unification handles tractable application patterns; unrestricted
higher-order unification is not assumed to be a cheap or complete subroutine.

To synthesise $h:\operatorname{List}(A)\to\operatorname{List}(B)$ with
$h([])=[]$ and $h(x::xs)=f(x)::h(xs)$, a list-fold schema already fixes the
recursion structure, so search needs only the empty case and the constructor
step; a library index may retrieve \code{map}, and otherwise a recursor term
provides an equivalent witness whose equations hold by reduction. Canonical is
the relevant prior work on inhabitation and unification in dependent type
theory, and it should be integrated through its proof-term interface rather than
by flattening away the dependencies that make its task
meaningful~\cite{canonical,canonical-min}.

After structural exposure, a saturation-based prover may find theorem
combinations a local tactic misses. Duper is an in-Lean proof-producing
option~\cite{duper}; \code{lean-auto} supplies translation and monomorphization
infrastructure~\cite{lean-auto}; \code{lean-smt} and QuerySMT supply SMT search
with Lean reconstruction~\cite{lean-smt,query-smt}. Two operational points:
\code{duper [*]} requests local facts, so a bare \code{duper} should not be
assumed to consume them automatically --- pass an explicit recorded premise
slice. And \code{lean-auto} distinguishes reconstruction-capable paths from
standalone invocation modes \emph{without} reconstruction, so an external
\code{unsat} is not a uniform proof interface.

The controller must not send every declaration in the local environment. Select
the dependency-relevant facts, instantiate polymorphic schemas only at the types
the current demand requires, and record which source theorem and which
universe/type instantiation each exported formula came from, so that
reconstruction applies the actual theorem.

A useful ordering for reconstruction, from most to least preferred: an existing
Lean tactic on the original expressions; a restricted, proved translation with a
known replay path; externally proposed hints replayed on the original
expressions; and finally a full external proof import with explicit supported
rules. \code{lean-smt}'s documented behaviour includes returning
proof-reconstruction holes as Lean goals; those holes enter the obligation graph
as children of the current AND node, and an external \code{unsat} response does
not bypass them. Partial progress is useful; it is not completed verification.

\subsection{Solver islands and the limits of theory combination}

A \emph{solver island} is a selected subproblem whose translation and proof
reconstruction are supported end to end. The island may prove a lemma over
shared interface terms rather than the complete original goal. For an island
under assumptions $A_1,\ldots,A_k$, reconstruct a proof of
$A_1\to\cdots\to A_k\to L$ and apply it to the appropriate scoped proofs to
insert $L$ into the main graph.

Every export carries a \emph{translation ledger}: the Lean expressions, their
solver symbols, their types and operations, and the proof obligations connecting
the encoding to the original proposition. Universe instantiation,
monomorphization, typeclass resolution, and auxiliary assertions must all have
justified transformations.

Abstraction has a direction, and both directions must be stated. Replacing a
difficult subterm by a fresh atom is sound for proving a \emph{universal
consequence} of the abstraction: a proof valid for every assignment of those
atoms is valid for the actual expressions. But an abstract satisfying assignment
can be spurious and is not a counterexample to the Lean goal. Conversely,
erasing dependencies in an indexed datatype is \emph{not} automatically a sound
abstraction; use an explicit bridge or leave that expression in the native
layer.

\begin{remark}[No blanket combination theorem]
Do not assume a generic Nelson--Oppen completeness result merely because solvers
exchange equalities. Finite datatypes and fixed-width bit-vectors do not satisfy
the usual stably-infinite assumptions. \forge{} avoids relying on a blanket
combination theorem: plugins exchange \emph{proved interface facts}. An
arithmetic atom abbreviating $xy$ must carry a definition or equality proof;
treating it as an unrelated real variable yields a relaxation, not a
bidirectional equivalence of problems.
\end{remark}

\subsection{Candidate generation can be untrusted without being unconstrained}

An optional model, learned ranker, or language model may propose a lemma, an
induction variable, a proof sketch, or a term. It cannot change the accepted
logical rules, and it cannot submit arbitrary Lean commands that introduce
axioms or alter the environment under the guise of a proof term. Restrict the
import channel to a typed certificate schema or to a term elaborated inside a
fixed environment, with resource and size limits protecting the checker from
pathological proposals. Never trust a textual theorem statement supplied by a
model.

Neural ranking is optional and is never proof evidence. The design does not
require an online model service: the executed prototypes use deterministic
enumeration, exact algebra, and local LP proposals, and their algorithms remain
useful in a fully offline implementation.

Finally, isolation is a separate concern from soundness. External search should
run in bounded child processes without network access --- not because a wrong
answer would be accepted, but because a foreign-function boundary is a
memory-safety surface and an unbounded search is a denial-of-service surface.
Neither is a soundness failure; both make a tactic unusable.
